\section{Additional Benchmark Details}
\label{sec:appendix-benchmark}

\subsection{Proposed Annotation Record}
\label{sec:annotation-record}

\paragraph{Source and event.}
Store the source identifier and URL, collection date, query provenance,
license or permission evidence, event-group identifier, and the original
start and end times of the evaluation unit. Record native audio,
embedded text, supplied subtitles, and unavailable content separately.
An event description should report what happens before cultural
interpretation is added.

\paragraph{Cue instance.}
Assign a stable identifier to each cue. Store its description, canonical
concept and aliases, component type, and one or more supported cultural
associations. Record the basis of each association and an uncertainty
flag. Evidence consists of channel labels and a set of time intervals;
optional image regions can support visually localized cues. A channel
should be marked only when it supports the specified cue, rather than
merely being present in the video.

\paragraph{Relation instance.}
Store the identifiers of participating cues, entities, or contexts, a
candidate relation frame, and the assignment of node identifiers to its roles.
Any cultural association attached to an entity or context needs a separate
evidence record. \todo{Confirm the node schema and treatment of integrated results.}
Record whether the relation is symmetric and whether the judgment is
supported by audiovisual evidence, additional metadata, or neither.
The ordering of node identifiers is not a semantic direction. Preserve
the evidence supporting the relation itself, since evidence that two
cues exist does not necessarily explain how they are connected.
\todo{Freeze field names and allowed values after the pilot.}

\subsection{Boundary Cases for Adjudication}
\label{sec:boundary-cases}
\begin{itemize}
    \item A recognizable practice occurs in a different country, but no
    cue in the video identifies the receiving context. This may require
    metadata and should not receive an unsupported audiovisual label.
    \item Two traditions share a practice. Several valid associations
    for that practice do not alone establish cultural mixing.
    \item Two forms appear successively in unrelated montage shots.
    Temporal proximity alone does not establish a shared event.
    \item A component is integrated into a base form, but the base cannot
    be identified reliably. Annotators should record the uncertainty
    rather than assign direction from the query order.
    \item A narrative adaptation is recognizable only from its title.
    The source--derivative relation belongs to the metadata-assisted
    condition unless the audiovisual evidence independently supports it.
\end{itemize}

\subsection{Protocol and Release Audit}
\label{sec:release-audit}
